\chapter{素数的阴谋}
\Large\textbf{The Prime Number Conspiracy: The Biggest Ideas in Math from Quanta}
\par \emph{Thomas Lin eds.} \normalsize

\par \textbf{集合论}：连续统假设(Continuum Hypothesis)被证明独立于ZFC公理系统（Gödel于1940年证明了ZFC无法推出CH为假, Cohen(1963)证明了ZFC无法推出CH为真）. 两个候选公理可被加入ZFC: (1)力迫公理(forcing axioms), 拓展集合的“宇宙”, 使连续统假设不成立, 可能有助于更多数学发现; (2)内模型公理(inner-model axiom) “V=终极L”, 即集合宇宙等于一个可构造、同时包含大基数的“内模型”（“较小”的宇宙）, 使连续统假设成立, 对数学其它领域影响较小. 

\par \textbf{代数}：McKay和Thompson发现了魔群不可约表示维数与j函数系数的联系. Conway-Norton(1979)提出的魔群月光(Monstrous Moonshine)猜想刻画了这一联系, 这一猜想被Borcherds(1992)借助弦论证明. 进一步地, 程之宁等(2014)提出的伴影月光(Umbral Moonshine)猜想断言了23种群与拟模形式系数的联系, 这23种新的“月光”分别有对应的弦论模型; 这一猜想被Gannon(2016), Duncan-Griffin-Ono(2015)证明. 

\par \textbf{数论}：张益唐(2014)证明了
\begin{equation*}
    \liminf\limits_{n\to \infty} (p_{n+1}-p_n)<7\times 10^7
\end{equation*}
这一结果改进了Goldston-Pintz-Yildirim(2009,2010)的结果，即相邻素数间隔与平均间隔的比例可以任意小：
\begin{equation*}
    \liminf\limits_{n\to \infty} \frac{p_{n+1}-p_n}{\log p_n}=0
\end{equation*}
Maynard(2015)将张益唐的上界改进到600，他的工作对素数三元组以至更多元组也给出了结果。

\par Bhargava(2004)给出了高斯二次型复合律的一种新表述，并发现了三次型的若干复合律。这一工作拓展了数的几何(geometry of numbers)，并被Bhargava用于解决数域与超椭圆曲线的问题。

\par Oliver-Soundararajan(2016)发现给定前一个素数的尾数，后一个素数尾数的分布不是均匀的，这可由Hardy-Littlewood(1923)的k元素数猜想解释。

\par \textbf{几何}：Mirzakhani的博士论文(2007,2008)解决了双曲曲面上简单测地线（不与自身相交的测地线）的计数问题，同时得出了模空间的体积公式，以及Witten猜想的一个新证明。

\par Scholze(2012)引入的拟完满空间(perfectoid spaces)改变了p进域下算术代数几何的面貌. 

\par 费曼图的振幅被发现与代数几何中的原相周期(period of motives)重合, 原相是各种上同调理论（奇异上同调，平展上同调，de Rham上同调）从不同角度刻画的同一主题。

\par \textbf{分析}：Hairer(2014)针对随机偏微分方程发展了正则结构(regularity structures)理论, 为物理学中的若干此类方程赋予了严格的数学含义。 

\par \textbf{概率论}：Tracy-Widom分布是一个左尾比右尾更陡的单峰分布, 不仅出现在一类随机矩阵最大特征值的分布中，还表示了弱耦合相（系统组分相对独立）与强耦合相（系统组分相互影响）间的普适相变。

\par Liouville量子引力(Liouville Quantum Gravity, LQG)和Brownian Map是随机二维曲面的两种描述模型，Miller-Sheffield(2020)利用Oded Schramm引入的非交叉随机曲线(Schramm-Loewner Evolution, SLE曲线)证明了二者的一致性。

\par \textbf{组合学}：Read猜想断言图的色多项式系数是单峰且对数凹的（$a_{n+1}a_{n-1}\le a_n^2$）；Rota猜想断言拟阵的特征多项式是对数凹的。许{\CJKfontspec{simsun.ttc}埈}珥 (2012)利用奇点理论证明了Read猜想，并与Adiprasito和Katz (2018)利用Hodge理论证明了Rota猜想。 

\par \textbf{离散几何}：Viazovska(2017)及合作者(2017)证明了8维、24维球体的最密堆积分别为$E_8$和Leech格。这一结果的基础是Cohen-Elkies(2003)提出用辅助函数估计球堆积密度上界；Viazovska从模形式中找到了合适的辅助函数。

\par 所有的三角形、凸四边形都可以密铺平面，可密铺的凸六边形有3种，而边数为7或更多的凸多边形不能密铺。Mann等(2018)找到了第15种可密铺的凸五边形；Rao(2017)用计算机验证了不存在其它的可密铺凸五边形。Penrose铺陈是由两种图形实现的非周期密铺；Socolar-Taylor(2011)给出了不连通单一图形非周期密铺的例子。Smith等(2024)给出了连通单一图形的非周期密铺（由于所有凸多边形的密铺都是周期密铺，这样的图形必是凹的）。与之相关的是密铺的可判定性问题，已证明一组图形的可密铺性是不可判定的，而单一图形的可判定性问题仍未知。

\par 王浩提出的“王氏砖”(Wang tiles)是一种边带颜色的正方形，只允许颜色匹配的边相拼接。已证实只存在非周期密铺的王氏砖的集合。Jeandel-Rao(2021)找到了一组只存在非周期密铺的11种王氏砖，并且11是最小的数字。

\par $\mathbb{F}_3^n$中不包含等差三元组的子集称为上集(cap set), 如SET游戏中不包含SET的一组不同纸牌对应$\mathbb{F}_3^4$中的一个上集, 其最大元素个数是20. Ellenberg-Gijswijt(2017)给出了上集个数的一个新上界. 